\documentclass[11pt]{article}
\usepackage[T1]{fontenc}
\usepackage[utf8]{inputenc}
\usepackage{lmodern,amsmath,amssymb,amsthm,mathtools}
\usepackage[margin=0.94in]{geometry}
\usepackage{microtype,booktabs,longtable,array,enumitem,xcolor}
\usepackage{hyperref}
\hypersetup{colorlinks=true,linkcolor=blue!45!black,urlcolor=blue!55!black,citecolor=blue!45!black,
pdftitle={Smaller explicit Diophantine forms of ZFC consistency},
pdfauthor={Research synthesis for the ProveIt project}}
\setlength{\emergencystretch}{3em}
\setlist{itemsep=3pt,topsep=5pt}
\newtheorem{theorem}{Theorem}[section]
\newtheorem{lemma}[theorem]{Lemma}
\newtheorem{proposition}[theorem]{Proposition}
\theoremstyle{remark}\newtheorem{remark}[theorem]{Remark}
\newcommand{\Z}{\mathbb Z}
\newcommand{\N}{\mathbb N}
\newcommand{\Pos}{\operatorname{Pos}}
\newcommand{\Con}{\operatorname{Con}}
\newcommand{\pc}{\operatorname{popcount}}
\newcommand{\Norm}{\operatorname{Norm}}
\newcommand{\KTM}{\mathsf{K215}}
\newcommand{\repo}{https://github.com/VladimirReshetnikov/ProveIt/blob/a477d0fc8b6feda53796835564a988015e744d4d/}
\newcommand{\coll}{SetTheory/Cardinals/docs/reports/hilbert-tenth-problem/con-zfc-diophantine-forms/}
\title{\textbf{Smaller explicit Diophantine forms\\of ZFC consistency}\\[5pt]
\large Twelve positive variables, twenty-four integers,\\and a smaller fixed coefficient tuple}
\author{Research synthesis prepared for the ProveIt project}
\date{22 October 2026}
\begin{document}
\maketitle
\begin{abstract}
We combine and refine the constructions in the audited ProveIt snapshot
\texttt{a477d0fc8b6f} (full pin in the references).
A new equation in twelve positive integers has exact degree \(1081\)
and a \(116\)-operation arithmetic circuit; its three-equation version
has maximum degree \(1072\) and \(110\) operations.
A different construction uses twenty-four unrestricted integers,
degree \(2432\), and \(163\) operations; its three-equation version
uses \(157\) operations and has maximum degree \(2416\).
A new parity layout reduces the largest fixed coefficient from
\(13{,}201{,}251\) to \(9{,}766{,}251\) bits while retaining the
943-tile computation language.
Ordinary mathematics formalizable in ZFC proves that each presentation
is solvable exactly when ZFC has a finite refutation, by composition
with the existing inputless compiler and exact machine calibration.
The archive supplies literal circuits, exact compressed coefficient
integers, finite recipes, proofs, independent reviews, and offline checks.
These are ordinary proofs, not new Lean or Rocq formalizations.
No global minimum or consistency proof is claimed.
\end{abstract}
\clearpage
\tableofcontents
\clearpage

\section{Results and the meaning of simplicity}

We seek a fixed polynomial \(P\) with integer coefficients such that
\[
\mathrm{ZFC}\vdash
\bigl(\neg\exists\mathbf x\in D^n\;P(\mathbf x)=0\bigr)
\longleftrightarrow\Con(\mathrm{ZFC}),
\tag{1.1}
\]
with an explicitly stated domain \(D\).
Here \(\N=\{0,1,2,\ldots\}\); positive integers are \(\Z_{>0}\).
Consistency means the absence of a finite syntactic refutation in the
recursive presentation used by the repository's exact machine endpoint.

There is no specified scalar measure of simplicity. We therefore keep
four measures separate: number of unknowns, degree, arithmetic circuit
length, and fixed coefficient size. A circuit is a literal finite
polynomial description: every gate uses integer addition, subtraction,
or multiplication. Shared expressions are abbreviations, not extra
existential variables.

\begin{table}[ht]
\centering\small
\begin{tabular}{llrrrr}
\toprule
New construction & Domain & Unknowns & Equations & Degree & Operations\\
\midrule
Relative root product & \(\Z_{>0}\) & 12 & 1 & 1081 & 116\\
Relative root system, shifted & \(\Z_{>0}\) & 12 & 3 & 1072 & 110\\
Signed saturation product & \(\Z\) & 24 & 1 & 2432 & 163\\
Signed saturation system, shifted & \(\Z\) & 24 & 3 & 2416 & 157\\
Signed saturation sum of squares & \(\Z\) & 24 & 1 & 4832 & 166\\
\bottomrule
\end{tabular}
\caption{For a system, degree means maximum residual degree.
Every binary operation is charged, including multiplication by a fixed
numeral; comparisons are free under the repository convention.}
\label{tab:new}
\end{table}

The positive product improves the previously recorded best degree among
the moderate-degree twelve-positive-variable equations, \(1145\) at
\(123\) operations. It matches the best \(116\)-operation count while
reducing that construction's degree from \(1353\).
The integer product improves the variable count, operation count,
and degree of every listed wave-3 twenty-six-integer product.
These are comparisons with the inspected repository family.

\subsection{The inspected baseline}

The audited reconciliation \cite{catalogue} records eight wave-3 packages,
numbered 16--23. Earlier packages survive in history and in reference
copies. The relevant recorded frontier is:

\begin{table}[ht]
\centering\small
\begin{tabular}{p{4.65cm}lrrr}
\toprule
Recorded construction & Domain & Unknowns & Degree & Operations\\
\midrule
Wave-1 \(P_{67}\), retained in 17/21 & \(\Z_{>0}\) & 15 & 155 & 67\\
Package 22, retained system & \(\Z_{>0}\) & 16 & 24 & 59\\
Package 20, lowest-degree root product & \(\Z_{>0}\) & 12 & 1145 & 123\\
Package 21, shortest root product & \(\Z_{>0}\) & 12 & 1353 & 116\\
Packages 21/22, root systems & \(\Z_{>0}\) & 12 & 1344/1520 & 112\\
Package 19, lower-degree integer product & \(\Z\) & 26 & 2465 & 192\\
Package 21, shorter integer product & \(\Z\) & 26 & 2657 & 168\\
Package 23, low-degree integer product & \(\Z\) & 27 & 172 & 116\\
Package 18, Jones specialization & \(\N\) & 9 & \(\sim1.9\cdot10^{28}\) & 20531\\
\bottomrule
\end{tabular}
\caption{Recorded baseline, not a fresh rebuild of every historical
package. The new work directly inspects the relevant proofs and
circuits in packages 19--23.}
\end{table}

Package 18's exact degree is \(161844\cdot5^{33}+484\).
It remains preferable when the sole objective is variable count in
that domain. The \(67\)-operation equation and the low-degree systems
also remain useful. We found no justified improvement below \(67\)
operations for a positive-domain single polynomial or below \(59\)
for the smallest recorded arithmetic system.

Package 20 supplies the direct root merger and rational inverse
formulas. Package 21 supplies an economical norm evaluator and Pell
saturation. Package 22 supplies the transport shift. Package 23
excludes a negative odd scale after reconstructing the core.
Package 19 supplies the Wang943 language. The new results combine
these compatible components through one precise interface.

\section{The inherited arithmetic and logical interface}
\label{sec:interface}

\subsection{Four fixed integer coefficients}

The coefficient ports are
\[
\beta=B-1,\qquad \kappa=D_C+BD_R,\qquad M_C,\qquad
\widehat M_F=M_F+B-1,\qquad B=2^{d_c}.
\tag{2.1}
\]
\(M_F\) is the native, unshifted high mask. The circuit port
\texttt{MF} is \(\widehat M_F\), and \texttt{K0} is \(\kappa\).
The four ports are fixed integers, not additional variables.
They satisfy
\[
\begin{gathered}
0<M_C,M_F<B-1,\quad M_C\equiv2\pmod4,\quad M_F\equiv4\pmod8,\\
\pc(M_C)+\pc(M_F)=d_c.
\end{gathered}
\tag{2.2}
\]
These conditions alone do not establish a compiler: the tile semantics,
alignment, carry bounds, and canonical completeness are also needed.
All four coefficients must be taken from one complete tuple.
Section~\ref{sec:coefficients} specifies the new exact tuple.

\subsection{Common polynomial abbreviations}

The twelve core coordinates are
\[
J,F,\alpha,z,h,i,T,s,w,\eta,\zeta,r.
\tag{2.3}
\]
Their domains will differ. Define
\begin{align}
q&=\beta J+1,& C&=q-F-\alpha,&X&=wq,&Y&=sq^3,&E&=XY,\nonumber\\
k&=\eta+\zeta,&c&=kY+\eta,&a&=E+Y,&H&=4a+3,&\Delta&=a^2+H,\tag{2.4}\\
\ell&=EkY,&A&=1+\ell(\ell+k),&
G&=1+\Delta c^2,&b_0&=X+ac.\nonumber
\end{align}
In particular \(\Delta=(a+2)^2-1\).
The rank expression and outer factors are
\begin{align}
R&=\bigl((q-1)(q-F)+\alpha\bigr)(q^2-1)
 +(M_C+q\widehat M_F)J,\tag{2.5}\\
N_k&=k-hE-R,\qquad
N_t=(\kappa+w)C+(q-F)-z(q-1).\tag{2.6}
\end{align}
The strong and auxiliary quantities are
\begin{align}
I&=ic^2,&D&=1+\Delta I^2,&Q&=\Delta(D-1)=(\Delta I)^2,&d&=Q-1,\nonumber\\
v&=cT,&b&=c+RD,&e&=Q(Dv^2+b^2)-1,&B_a&=2Qvb.
\tag{2.7}
\end{align}
Every letter in these displays is a polynomial abbreviation.
All exponents are fixed nonnegative integers.

\subsection{The reused core theorem}

The inputless proofs \cite{blank,ordinary,signed,oddscale} give an
interface with reconstructed integer roots \(\tau,u,f,y\) and
integer \(\gamma\):
\begin{align}
\tau^2-\ell(\ell+k)&=1,\nonumber\\
u=b_0+H\gamma,\qquad u^2-\Delta c^2&=1,\nonumber\\
f^2-\Delta i^2c^4&=1,\tag{2.8}\\
Q[c(Tf-1)-Rf^2]^2-(Q-1)y^2&=1.\nonumber
\end{align}
When \(J,F,\alpha,z,s,w,\eta,\zeta>0\) and \(i\ne0\), the other
coordinates \(h,T,\gamma,f,\tau,y\) may be signed.
Together with \(N_k=N_t\in\{1,-1\}\), these equations force
both outer factors to \(+1\) and decode a blank halt of \(\KTM\).
The ordinary-strong interface is obtained using
\(i_{\rm ordinary}=\Delta|i|\), leaving \(Q\) unchanged.

The proof initially gets only \(C\ge0\), not \(C>0\). Its bounds are
\[
q\ge16,\quad q(q^2-1)<R<q^4-2,\quad R-2>3q+1.
\tag{2.9}
\]
With positive scales \(E\ge q^4>R+2\).
The source Pell-rank argument recovers the exact main rank \(R\)
before settling the unit signs and extracting the tableau.
This proof order preserves the exceptional \(C=0\) boundary correctly.

Conversely, a finite blank halt has canonical positive witnesses.
Independent width/height padding by powers of five can simultaneously
give odd \(i\), square \(w\), and odd \(s\). The first two choices
use the normalized strong index \(m=cR\) and an odd exponent of \(q\);
the exact population calculation for \(s\) is recalled below.
These choices preserve existence and are not claimed for every
arbitrary earlier witness.

\subsection{The ZFC endpoint and proof status}

The inspected Lean endpoint \cite{endpoint} contains
\[
\operatorname{BlankHalts}(\KTM)
\longleftrightarrow\operatorname{BProv}(\mathrm{ZFCax}_s,\bot),
\tag{2.10}
\]
named \texttt{KEndpoint215.blank\_halts\_iff\_zfc\_refutation}.
Its nonhalting theorem is obtained by negation.
The ordinary finite-computation and syntactic-refutation calibration
is the logical input here. We do not infer internal ZFC provability
from an assertion that Lean is sound.
The compiler and the new reductions are ordinary mathematics
internalizable in ZFC; their composition proves (1.1).

No endpoint build was rerun, and the new machine-to-Diophantine
composition is not a newly checked Lean or Rocq theorem.
The finite audits certify their stated finite properties, while the
mathematical arguments establish the unbounded equivalences.

\section{A polynomial norm without hidden witnesses}
\label{sec:norm}

Reduce \(p(t)\) modulo the monic polynomial \(t^2-g\) to \(p_0+p_1t\)
and define
\[
\Norm_{g,t}(p)=p_0^2-gp_1^2.
\tag{3.1}
\]
This is an explicit finite operation on polynomials. It adds no
existential root and requires no root extraction at evaluation time.
For polynomial weights \(L,M\) and \(W=b_0+Hr\), put
\[
\begin{split}
\mathcal N(L,M,W)=\Norm_{G,u}\Norm_{A,t}\bigl(
 &[d(W-Lu-HMt)^2-H^2e]^2\\
 &-H^4B_a^2D\bigr).
\end{split}
\tag{3.2}
\]
Equivalently this is the eight-factor product
\[
\prod_{\substack{u=\pm\sqrt G\\t=\pm\sqrt A\\f=\pm\sqrt D}}
\bigl(d(W-Lu-HMt)^2-H^2e+H^2B_af\bigr).
\tag{3.3}
\]
All branches are counted with multiplicity, including coincident
quadratic extensions.

\subsection{A literal evaluator}

The following expressions expand the norm using only \(+,-,\times\):
\begin{align}
Z_0&=W^2+L^2G,&Z_1&=2WL,&n&=(HM)^2A,\nonumber\\
C_0&=d(Z_0+n)-H^2e,&C_1&=dZ_1,&K_f&=4d^2n,\nonumber\\
K'_0&=K_fZ_0,&K'_1&=K_fZ_1,&j&=H^4B_a^2D,\tag{3.4}\\
V_0&=C_0^2+GC_1^2-K'_0-j,&V_1&=2C_0C_1-K'_1,\nonumber\\
N_0&=V_0^2+GV_1^2-4jK'_0,&N_1&=2V_0V_1-4jK'_1,\nonumber\\
\mathcal N&=N_0^2-GN_1^2.\nonumber
\end{align}
The scalar \(j\) is evaluated as \(4H^4PU\), where \(P=QDv^2\)
and \(U=Qb^2\). The identity
\[
(C^2+j-K)^2-4jC^2=(C^2-K-j)^2-4jK
\tag{3.5}
\]
in the quadratic ring explains the economical evaluator.
\(K'_0\) is unrelated to the fixed coefficient \(\kappa\).

\begin{lemma}[Unit guard]\label{lem:unit}
For any integer assignment, \(8\mid N_1\), so
\(\mathcal N\equiv N_0^2\pmod{64}\). Therefore
\[
N_kN_t(1+\mathcal N)=1
\ \Longrightarrow\ 
\mathcal N=0,\quad N_k=N_t\in\{1,-1\}.
\tag{3.6}
\]
\end{lemma}
\begin{proof}
\(Z_1,C_1\) are even, \(2C_0C_1\) is divisible by four,
and \(K'_1\) is divisible by eight. Hence \(4\mid V_1\)
and \(8\mid N_1\). At a product zero all three factors are integer
units. If \(1+\mathcal N=-1\), then \(\mathcal N=-2\equiv2\pmod4\),
contrary to its square residue. Thus the last factor is \(1\).
\end{proof}

\section{Rational Pell saturation}
\label{sec:saturation}

\begin{lemma}\label{lem:saturation}
If \(m\ge2\) and \(V\) are integers and \(y\in\mathbb Q\) satisfies
\[
m^2V^2-(m^2-1)y^2=1,
\tag{4.1}
\]
then \(y\) is an integer.
\end{lemma}
\begin{proof}
Write \(y=A_0/B_0\) in lowest terms with \(B_0>0\).
Equation (4.1) gives \(B_0^2\mid m^2-1\).
The positive integer \(d_0=(m^2-1)/B_0^2\) is nonsquare.
The two positive integer-coordinate Pell solutions are
\((m,B_0)\) and \((m|V|,|A_0|)\); their ordinates are nonzero.
Write the powers of the least positive integer-coordinate norm-one
unit as \((\chi_j,\psi_j)\). For positive \(j,l\),
\[
m=\chi_j,\ B_0=\psi_j,\ m|V|=\chi_l,\ |A_0|=\psi_l.
\]
The recurrence gives
\(\chi_{n+2j}\equiv-\chi_n\pmod{\chi_j}\) and
\(\chi_{2j-n}\equiv-\chi_n\pmod{\chi_j}\).
Reduce \(l\) modulo \(2j\) and reflect to \(0\le n\le j\).
Since \(0<\chi_n<\chi_j\) for \(n<j\), divisibility
\(\chi_j\mid\chi_l\) forces \(l\) to be an odd multiple of \(j\).
Pell multiplication gives \(\psi_j\mid\psi_l\), hence \(B_0\mid A_0\);
coprimality forces \(B_0=1\).
\end{proof}

This saturation lemma is inherited from packages 19--22
\cite{norm21,root20,monic22}. It closes an integrality gap that
rationality alone would not close for a general nonmonic quadratic.
The new work combines it with shorter weights and with a signed
scale that removes two positive-domain witnesses.

\section{A smaller equation in twelve positive integers}
\label{sec:positive}

Use all twelve coordinates in (2.3) as positive integers. Define
\[
M=kI^2+T^2,\qquad L=1+2HM,\qquad W=b_0+Hr.
\tag{5.1}
\]
The proposed single equation is
\[
\boxed{\quad P_{12}:=N_kN_t\bigl(1+\mathcal N(L,M,W)\bigr)-1=0.\quad}
\tag{5.2}
\]
Together with the finite definitions in Sections~\ref{sec:interface}
and \ref{sec:norm}, this is a fixed integer polynomial, not an
exponential Diophantine equation.

\begin{theorem}\label{thm:positive}
For the complete compiler tuple of Section~\ref{sec:coefficients},
\(P_{12}=0\) has a solution in twelve positive integers if and only if
\(\KTM\) halts on its blank tape. Its exact total degree is \(1081\);
the supplied circuit uses \(116\) binary arithmetic operations.
\end{theorem}

\subsection{Weak bounds before root recovery}

By Lemma~\ref{lem:unit}, a product zero gives
\(\mathcal N=0\) and \(N_k=N_t=\epsilon\in\{1,-1\}\).
The inherited transport bootstrap is unchanged, since it involves only
the outer expressions. It gives
\[
\begin{gathered}
q\ge16,\quad X\ge16,\quad Y\ge4096,\quad C\ge0,\quad
0<R<q^4-2,\quad R>3q\ge48,\\
E\ge q^4>R+2,\qquad
k=R+\epsilon+hE>E,\qquad R<k,\\
kY<c<k(Y+1)<2kY.
\end{gathered}
\tag{5.3}
\]
No square-root recovery, sign of \(\gamma\), or strict inequality
\(C>0\) is assumed in this step.

\subsection{The inner weight}

Write \(a_0=a+2\) temporarily. Since \(X\ge16\), the weak bounds give
\(a_0<2E\). They also give
\[
\ell=kEY>\frac{a_0c}{4},\qquad
\ell>E^2Y>8a_0^2,
\]
and therefore
\[
A>\frac{a_0^2c^2}{16},\qquad A>64a_0^4.
\tag{5.4}
\]
As \(I=ic^2\) with \(i\ge1\), we have
\[
D<a_0^2I^2,\qquad c<I^2,\qquad b=c+RD<2ka_0^2I^2.
\]
Using \(M^2\ge4kI^2T^2\) and \(k>48\), we obtain two separate estimates:
\[
AM^2>\frac{k}{4}Dv^2>8Dv^2,\qquad
AM^2\ge Ak^2I^4>16b^2>8b^2.
\]
Averaging them, and using \(d=Q-1\ge Q/2\), yields
\[
\boxed{\quad dAM^2>2Q(Dv^2+b^2)=2(e+1).\quad}
\tag{5.5}
\]
This strengthening of the inherited constants allows us to use
\(M=kI^2+T^2\) without an extra factor of two.

For either root \(f'=\pm\sqrt D\), any complex \(y'\) satisfying
\(d(y')^2=e-B_af'\) obeys
\[
\begin{split}
d|y'|^2
&\le e+B_a\sqrt D
 <Q(v\sqrt D+b)^2\\
&\le2Q(Dv^2+b^2)<dAM^2.
\end{split}
\]
Thus every such branch satisfies
\[
|y'|<M\sqrt A.
\tag{5.6}
\]
The complex absolute value matters: the right-hand side of the
quadratic equation for a particular branch can be negative.

\subsection{Relative separation of the two main-root branches}

The decisive simplification is the comparison \(A<G\):
\[
\begin{split}
G-A
&>k^2\bigl((E+Y)^2Y^2-EY(EY+1)\bigr)\\
&=k^2\bigl(EY(2Y^2-1)+Y^4\bigr)>0.
\end{split}
\tag{5.7}
\]
Let \(\Phi(x)\in\Z[x]\) be the product over
\(t=\pm\sqrt A,\ f=\pm\sqrt D\) of
\[
d(x-Mt)^2-e+B_af.
\]
Every complex root of \(\Phi\) has the form \(x=Mt+y'\), so (5.6)
gives the uniform bound
\[
|x|<2M\sqrt A<2M\sqrt G.
\tag{5.8}
\]
At \(\mathcal N=0\), one root \(u=\pm\sqrt G\) gives
\(x=(W-Lu)/H\) with \(\Phi(x)=0\).
If \(G\) is not a square, conjugation in \(\mathbb Q(\sqrt G)\)
also makes \(x'=(W+Lu)/H\) a root. But
\[
|x'-x|=\frac{2L\sqrt G}{H}>4M\sqrt G,
\]
contradicting the sum of the two bounds (5.8).
Consequently \(G\) is a square and \(u\) is an integer.
Moreover \(u>0\): if \(u<0\), then
\[
W=Lu+Hx<-L\sqrt G+2HM\sqrt G=-\sqrt G<0,
\]
whereas \(W=b_0+Hr>0\).
We have recovered an integer main root and a rational \(x\).

The earlier absolute separation used an additional factor
\(\ell+k\) in \(L\). Comparing the inner radius with the distance
between the two main-root branches removes that factor.

\subsection{Recovering the remaining integer roots}

For this rational \(x\), put
\[
\mathcal C=d(x^2+M^2A)-e,\qquad
\mathcal V=\mathcal C^2+DB_a^2-4d^2x^2M^2A.
\tag{5.9}
\]
Equation (5.5) implies \(\mathcal C>e+2>0\).
Also \(x\ne0\), since
\[
\Phi(0)=\bigl[(dM^2A-e)^2-DB_a^2\bigr]^2>0,
\]
using \(B_a\sqrt D\le e+1\).
The coherent inverse formulas inherited from the direct merger are
\[
f=-\frac{\mathcal V}{2\mathcal C B_a},\qquad
t=\frac{\mathcal C+B_af}{2dxM},\qquad
y=x-Mt.
\tag{5.10}
\]
Their denominators are nonzero. Squaring the original factor and
eliminating \(t^2=A\) gives
\(\mathcal V+2\mathcal C B_af=0\); substituting back gives
the displayed expression for \(t\).
Thus (5.10) supplies rational numbers with
\(f^2=D\), \(t^2=A\), and
\[
Q(vf-b)^2-(Q-1)y^2=1.
\tag{5.11}
\]
Rational square roots of integers are integers, so \(f,t\in\Z\).
Since \(Q=(\Delta I)^2\) and \(vf-b\in\Z\),
Lemma~\ref{lem:saturation} implies \(y\in\Z\).

Now \(x=Mt+y\) is integral. Reducing \(W=Lu+Hx\) modulo \(H\),
and using \(L\equiv1\pmod H\), gives
\[
\gamma=\frac{u-b_0}{H}\in\Z.
\]
In fact \(\gamma>0\): from \(c>k>E>X\),
\(\Delta>(a+1)^2\), and \(u>0\), we get
\(u>(a+1)c\), hence \(\gamma H>c-X>0\).
The identity
\[
vf-b=c(Tf-1)-Rf^2
\]
uses \(f^2=D\) and restores all four equations (2.8).
Their permitted signed quotient absorbs a negative \(f\) by changing
\((f,T)\) to \((-f,-T)\).
The inherited core theorem forces \(\epsilon=1\) and decodes a
blank halt of \(\KTM\).

Conversely, take a canonical positive normalized compiler witness.
Choose its positive roots \(u,t,y\) and set
\[
r=\gamma+2Mu+Mt+y>0.
\tag{5.12}
\]
Then \(W=Lu+H(Mt+y)\), so one factor in (3.3) vanishes.
The outer factors both equal one. This proves both directions of
Theorem~\ref{thm:positive}.

\subsection{A three-equation version in 110 operations}

Initially the equivalent system is
\[
N_k=1,\qquad N_t=1,\qquad N_0^2=GN_1^2.
\tag{5.13}
\]
Direct comparison in the last equation saves the subtraction
defining \(\mathcal N\).
For the shorter system replace the old transport coordinate by
\(z_{\rm old}=z+1\) and use
\[
\boxed{
\begin{aligned}
N_k&=1,\\
(\kappa+w)C&=F+z(q-1),\\
N_0^2&=GN_1^2.
\end{aligned}}
\tag{5.14}
\]
Every positive solution of (5.14) gives one of (5.13).
Conversely, after the compiler has been recovered, it gives
\(C>0\) and \(X=wq=2^R\) with \(R>3q\). Hence \(w>q>F\), and
\[
(\kappa+w)C-F=(z_{\rm old}-1)(q-1)>0.
\]
Thus \(z_{\rm old}\ge2\), making \(z=z_{\rm old}-1\) positive.
This is the zero-set shift from package 22 \cite{monic22};
it is used after soundness, not as a premature weak bound.
The product (5.2) retains the old coordinate.

The product ledger is \(74\) multiplications and \(42\) additions
or subtractions. The shifted system uses \(72+38=110\) operations;
its residual degrees are \(7,2,1072\).

\subsection{Exact degree certificate}

Give all twelve supplied variables degree one and the four ports
degree zero. The relevant degrees are
\[
\begin{array}{c|rrrrrrrrr}
 &c&H&\Delta&I&D&Q&R&b&e\\ \hline
\deg&5&6&12&11&34&46&4&38&122
\end{array}
\]
and
\[
\begin{array}{c|rrrrrrr}
 &A,G&M&L& W&n&Z_0&Z_1\\ \hline
\deg&22&23&29&11&80&80&40
\end{array}
\]
The second array has a shared \(A,G\) entry.
In \(C_0\) the unique highest homogeneous part is \(-H^2e\),
of degree \(134\), while \(d(Z_0+n)\) has degree at most \(126\).
It follows successively that \(V_0\) has highest part \(H^4e^2\),
of degree \(268\), and \(N_0\) has highest part \(H^8e^4\),
of degree \(536\). In contrast, \(\deg N_1\le488\).
Therefore the final terms have degrees \(1072\) and at most \(998\);
their leading terms cannot cancel.

More explicitly, with \(q_0=\beta J,\ k_0=\eta+\zeta\),
\[
\mathcal N_{\rm top}
=2^{32}i^{48}k_0^{96}w^{80}s^{176}q_0^{656}(q_0-F)^{16}.
\tag{5.15}
\]
The product's highest part is this times
\[
(-hwsq_0^4)\,[w(q_0-F-\alpha)-zq_0].
\]
The coefficient of
\[
J^{677}hzi^{48}\eta^{96}w^{81}s^{177}
\]
is \(2^{32}\beta^{677}\ne0\).
Thus \(\deg\mathcal N=1072\) and \(\deg P_{12}=1081\)
for every authentic tuple, not just a generic small specialization.

\section{Twenty-four unrestricted integer unknowns}
\label{sec:integer}

The positive proof above does not extend merely by allowing \(s\)
to be negative: its comparison \(A<G\) and the signs of its weights
were proved on the positive domain. For the integer construction we
instead retain the polynomial norm circuit of package 21 and prove a
larger domain for it.

Use the common abbreviations (2.4)--(2.7), but now set
\[
\mathcal B=v^2+D+2b,\qquad
L=1+H\mathcal B,\qquad M=L,\qquad W=b_0+Hr.
\tag{6.1}
\]
Require positivity only for
\[
J,F,\alpha,z,\eta,\zeta.
\]
For each of these six coordinates use a distinct triple of integer
unknowns and the polynomial
\[
\Pos(x,y,t)=1+x^2+y^2+t(t+1).
\tag{6.2}
\]
It is positive for every integer triple and maps onto the positive
integers. Indeed \(4n-3\equiv1\pmod4\) is a sum of three squares
with one odd square and two even squares; dividing after separating
the odd square gives (6.2).

For the remaining coordinates put
\[
i=2i_0+1,\qquad s=2s_0+1,\qquad w=w_0^2,
\qquad h,T,r\in\Z.
\tag{6.3}
\]
There are \(6\cdot3+6=24\) independent integer unknowns.
After these substitutions the polynomial is
\[
\boxed{\quad P_{24}:=N_kN_t
 \bigl(1+\mathcal N(L,L,W)\bigr)-1=0.\quad}
\tag{6.4}
\]

\begin{theorem}\label{thm:integer}
For the complete compiler tuple of Section~\ref{sec:coefficients},
\(P_{24}=0\) has an integer solution if and only if \(\KTM\)
halts on its blank tape. Its exact degree is \(2432\), and its
circuit uses \(163\) binary operations.
\end{theorem}

\subsection{Bounds valid for both signs of the scale}

The unit guard again gives \(\mathcal N=0\) and
\(N_k=N_t=\epsilon\in\{1,-1\}\). Positivity of \(F,z\) implies
\[
(q-F)-z(q-1)=1-F-(z-1)(q-1)\le0.
\]
Since \(\kappa+w\ge2\), the alternative \(C\le-1\) would make
\(N_t\le-2\). Hence \(C\ge0\), and the unchanged mask argument gives
the bounds (2.9).

Put \(Z=|Y|\). The odd substitution excludes \(s=0\), so
\[
\begin{gathered}
Z\ge q^3,\quad X\ge0,\quad |a|=(X+1)Z,\quad
k(Z-1)<|c|<2kZ,\quad |H|>1,\\
\Delta\ge(q^3-1)(q^3-3)>q^4>R.
\end{gathered}
\tag{6.5}
\]
Since \(i\ne0\) and \(|c|>1\),
\[
D>\Delta+|c|,\qquad
1<b=c+RD<Q.
\tag{6.6}
\]
For the upper bound, explicitly,
\[
Q-b=(\Delta-R)D-\Delta-c
 \ge D-\Delta-|c|>0.
\]
Also \(G,D,Q>1\) and
\(\ell=XkY^2\ge0\), so \(A\ge1\).
In particular \(\mathcal B>0\), even with signed or zero \(T\).

\subsection{Root recovery with possibly negative weights}

The factorization of the norm supplies a coherent triple
\(u^2=G,\ t^2=A,\ f^2=D\) and
\[
y=\frac{W-Lu-HLt}{H},\qquad
dy^2=e-B_af=Q(vf-b)^2-1.
\tag{6.7}
\]
Work in a finite normal splitting field. Every conjugate \(y'\)
satisfies \(|y'|<\mathcal B\).
To see this, put \(M_0=|v|\sqrt D+b>1\).
The real quantity \(Q(vf'-b)^2-1\) lies between \(-1\) and
\(QM_0^2-1\), so its absolute value is at most \(QM_0^2-1\).
Thus
\[
|y'|^2\le\frac{QM_0^2-1}{Q-1}<2M_0^2,\qquad
\mathcal B=v^2+D+2b\ge2M_0.
\tag{6.8}
\]
This bounds imaginary as well as real branches.

Set \(U_0=|H|\mathcal B>2\).
Although \(L\) can be negative, \(|L|\ge U_0-1\).
Since \(\sqrt G>2\), we obtain
\[
|L|\sqrt G>U_0.
\tag{6.9}
\]
If \(X>0\), then \(X\ge q\), and
\[
\sqrt A>\ell=kZ|E|,\qquad
\sqrt G<(|a|+2)|c|<2(|a|+2)kZ.
\]
Consequently
\[
\begin{split}
|H|\sqrt A-\sqrt G
&>kZ\bigl((4|a|-3)|E|-2(|a|+2)\bigr)\\
&\ge kZ(2|a|-7)>2.
\end{split}
\]
Using \(2|L|>U_0\), this implies the second strict separation
\[
|HL|\sqrt A>|L|\sqrt G+|H|\mathcal B.
\tag{6.10}
\]
If \(t\) were irrational, an automorphism of the splitting field
would flip its sign. Subtracting the conjugates of
\(W=Lu+HLt+Hy\) would contradict (6.10), because the other two
differences have absolute values at most \(2|L|\sqrt G\)
and strictly less than \(2|H|\mathcal B\).
Thus \(t\) is rational.
If \(X=0\), then \(A=1\), so \(t\) is rational directly
and (6.10) is unnecessary.

With \(t\) rational, an automorphism flipping \(u\) would contradict
(6.9). Hence \(u\) is rational as well. Both are integers because
their squares are integers; equation (6.7) then makes \(y\) rational.

\subsection{Saturation and the two exceptional boundaries}

If \(T\ne0\), then \(v\ne0\), \(B_a\ne0\), and
\[
f=\frac{e-dy^2}{B_a}\in\mathbb Q.
\]
Since \(f^2=D\), this root is integral. Now \(vf-b\in\Z\), so
Lemma~\ref{lem:saturation} applied with
\(m=|\Delta ic^2|\) makes \(y\) integral.
If \(T=0\), the same lemma applies directly with \(V=b\),
without first recovering \(f\), and again gives \(y\in\Z\).
Because \(L\equiv1\pmod H\), in either case
\[
\gamma=\frac{u-b_0}{H}
=r-\mathcal B u-Lt-y\in\Z.
\tag{6.11}
\]

If \(X=0\), then \(u=ac+H\gamma\).
Reducing \(u^2-\Delta c^2=1\) modulo \(H\), with
\(\Delta=a^2+H\), gives \(H\mid1\), contrary to \(|H|>1\).
Therefore \(X>0\) and \(w>0\).
If \(T=0\), (6.7) gives
\[
y^2=b^2+\frac{b^2-1}{Q-1}.
\]
The integer inequalities \(1<b<Q\) put this number strictly
between \(b^2\) and \((b+1)^2\), a contradiction.
Thus \(T\ne0\), and \(f\) is integral in all possible cases.

We have now reconstructed all four equations (2.8), with signed
\(t,u,f,y,\gamma,h,T\). No division by an unproved nonzero
\(T\) or \(B_a\) was used.

\subsection{Why a negative odd scale is impossible}

The recovered equations meet the hypotheses of the signed-scale
exclusion in package 23 \cite{oddscale}. The normalized multiplier
need only be nonzero: replacing \(i\) by \(|i|\) changes none of
the equations.
For clarity, the mechanism of that exclusion is as follows.
If \(s<0\), put
\[
Z=-Y,\qquad B_0=(X+1)Z-2,\qquad c_0=-c,\qquad
P_0=2XZ^2+1.
\]
The first and main Pell classifications give indices \(n,p\)
with \(k=2\psi_{P_0}(n)\) and \(c_0=\psi_{B_0}(p)\).
Here \(\psi_B\) denotes the ordinate of
\((B+\sqrt{B^2-1})^j\).
The normalized strong equation makes its index divisible by
\(pc_0\). The auxiliary congruences then recover \(p=R\);
the sign of \(c\) disappears in the squared congruence.
The strict index interval forces \(\epsilon=1\) and
\(R=2r_0+1\), and the Pell growth comparison gives
\[
Z>\frac{X^{r_0}}2>2^R.
\]
The two possible signs of the main root give respectively
\[
X+2^R\equiv0\pmod H
\quad\text{or}\quad
X\,2^R\equiv1\pmod H.
\]
Both are excluded by
\[
0<X+2^R<|H|,
\qquad
1<X\,2^R<XZ<|H|.
\]
The full inherited rank argument, including its signed congruences,
is supplied in the reference copy of
\texttt{ODD\_SCALE\_27\_PROOF.md}.
It applies to the reconstructed core, not to a merely rational
norm witness. This order is essential.

Since \(s\) is odd and cannot be negative, it is positive.
The signed normalized compiler now decodes a blank \(\KTM\) halt
and forces \(N_k=N_t=1\).

\subsection{Simultaneous completeness of the domain restrictions}

From an actual blank halt, canonical independent width and height
padding gives \(q=2^\theta\) with \(\theta=d_cN\) odd.
The canonical rank \(R\) is odd, so
\[
w=2^{R-\theta}
\]
is a positive square.
Choosing the normalized strong index \(m=cR\) gives odd positive
\(i\), together with the required auxiliary quotient, as in
the signed-domain proof \cite{signed}.
These restrictions concern the same padded tableau.

To check odd \(s\) without changing that witness, set
\[
S=(C-1)+qF,\qquad
T_{\rm mask}=(M_CJ+1)+q(M_FJ-1),
\tag{6.12}
\]
where \(M_F\) is the \emph{native} high mask, not the shifted
circuit port \(\widehat M_F\).
The native mask constraints give
\(S\mathbin{\&}T_{\rm mask}=0\), both quantities below \(q^2\),
and \(\pc(T_{\rm mask})=\theta+2\).
The exact packing identity is
\[
R=(q^2-1-S)q^2+(S+T_{\rm mask}).
\]
The two blocks are disjoint, so
\[
\pc(R)=2\theta+\pc(T_{\rm mask})=3\theta+2.
\]
The inherited half-binomial valuation gives
\(v_2(Y)=\pc(R)-2=3\theta\). Therefore \(s=Y/q^3\) is positive odd.
Every remaining guarded coordinate is in the image of (6.2).
Finally, with \(t\) again the canonical first Pell root, choose
\[
r=\gamma+\mathcal B u+Lt+y.
\tag{6.13}
\]
This integer makes one norm factor vanish, proving completeness
for (6.4) with all three restrictions satisfied simultaneously.

\subsection{The system and sum-of-squares alternatives}

The unshifted system is \(N_k=1,\ N_t=1,\ \mathcal N=0\).
For the recommended shorter system, use six positive maps including
the new \(z\), put \(z_{\rm old}=z+1\), and write
\[
\boxed{
\begin{aligned}
N_k&=1,\\
(\kappa+w)C&=F+z(q-1),\\
N_0^2&=GN_1^2.
\end{aligned}}
\tag{6.14}
\]
The proof of the shift is the same as for (5.14), after the core
has been decoded. This system has \(24\) integer unknowns,
three equations, \(157\) gates, and residual degrees \(12,4,2416\).
The unshifted version uses \(159\) gates.
For a single polynomial obtained directly from the system, the
archive also includes
\[
(N_k-1)^2+(N_t-1)^2+\mathcal N^2=0,
\tag{6.15}
\]
with degree \(4832\) and \(166\) operations.
Its integer zero set is exactly that of the unshifted system.

\subsection{Operation count and exact degree}

Each positive map costs seven operations, the two odd maps cost two
each, and the square map costs one. The total domain-conversion
cost is therefore \(42+4+1=47\).
The \(116\)-gate package-21 product circuit is retained literally
after that prefix, yielding \(163\) gates.
This is a new theorem about the domain of an old polynomial
circuit, as well as a smaller explicit integer presentation.

Give \(J,F,\alpha,z,w,\eta,\zeta\) weight two and the other
core coordinates weight one. Key weighted degrees are
\[
\begin{array}{c|rrrrrrrr}
 &c&H&\Delta&A,G&D&Q&b,\mathcal B&L\\ \hline
\deg&9&11&22&40&60&82&68&79
\end{array}
\]
In the evaluator (3.4), \(n\) has weight \(220\), and the unique
highest part of \(C_0\) is \(dn\), of weight \(302\).
The highest part of \(N_0\) has weight \(1208\), while
\(\deg N_1\le1087\).
Thus \(N_0^2\) has weight \(2416\) and
\(GN_1^2\) has weight at most \(2214\); there is no leading
cancellation. With \(q_0=\beta J,\ k_0=\eta+\zeta\), the
highest weighted part of the norm is
\[
2^{80}i^{48}k_0^{112}w^{112}s^{224}q_0^{832}(q_0-F)^{16}.
\tag{6.16}
\]
Multiplying by the leading outer factors gives
\[
\begin{split}
(P_{24})_{\rm top}={}&-2^{80}h i^{48}k_0^{112}w^{113}s^{225}
 q_0^{836}(q_0-F)^{16}\\
&\hspace{9mm}\cdot[w(q_0-F-\alpha)-zq_0],
\end{split}
\tag{6.17}
\]
before the actual degree-two substitutions.
For a guarded coordinate \(p\), denote the first variable in its
positive-map triple by \(p_x\).
After substituting, the coefficient of
\[
J_x^{1706}h z_x^2 i_0^{48}\eta_x^{224}w_0^{226}s_0^{225}
\]
is \(2^{353}\beta^{853}\ne0\).
Its degree is \(2432\).
The analogous norm coefficient is
\(2^{352}\beta^{848}\ne0\), at degree \(2416\).
Consequently the product and system degrees hold for the actual
fixed coefficients, uniformly in every authentic \(\beta>0\).

\section{A smaller complete coefficient tuple}
\label{sec:coefficients}

The preceding reductions can use a complete authentic tuple from
the inherited compiler. We additionally give a smaller tuple for
the same \(943\)-tile Wang language.
The improvement changes the arithmetic placement of tile-copy
tests, while retaining the ordered language, the blank-machine
semantics, and the radix-cell size.

\subsection{Exact finite specification}

The archive includes the ordered list of \(943\) tiles from
package 19 \cite{wang}, generated from the literal \(\KTM\)
transition table. A tile is
\((x_W,x_N,x_E,x_S)\), with all labels below \(512\).
Selector \(0\) is Start; selector \(1\) is the globally impossible
Forbidden tile. The list's SHA-256 is
\[
\begin{gathered}
\texttt{66b55f1d11fca9fefc93b12fcfc8b4cabb75b}\\
\texttt{08e9fd03311f9e894a1c5e5e81f}.
\end{gathered}
\]
The split display denotes one \(64\)-digit hash with no separator.
This fixed finite list, rather than a hypothetical halting run,
is the input to the following recipe.

Set
\[
A_c=1024,\quad \rho=2^{625},\quad L_c=15625,\quad
d_c=625L_c=9765625=5^{10},\quad B=\rho^{L_c}.
\tag{7.1}
\]
For \(j=0,\ldots,9\), place the four kinds of native copy bits at
\[
\begin{array}{c|rrrr}
t&W&N&E&S\\ \hline
e_t(j)&1020+2j&1000+2j&980+2j&960+2j.
\end{array}
\tag{7.2}
\]
Let \(\mathcal E\) be the disjoint union of selectors \(0,\ldots,942\),
these forty copy positions, dummy positions \(1063,1064\), and
anchors \(1095,1096,2195,3296\).
Thus \(|\mathcal E|=989\), and \(\max\mathcal E=3296\).
Order directions as \(W,N,E,S\), numbered \(t=0,1,2,3\).
For selector \(i\), define the fixed integer
\[
c_i=1+\sum_{t=0}^3\sum_{j=0}^9
 \operatorname{bit}_j(x_t+512)A_c^{1+10t+j}
 +\sum_{j=41}^{44}A_c^j.
\tag{7.3}
\]
The remaining coefficients are
\[
\begin{gathered}
c_{e_t(j)}=A_c^{1+10t+j},\qquad
c_{1095}=A_c^{41},\quad c_{1096}=A_c^{42},\\
c_{2195}=A_c^{43},\quad c_{3296}=A_c^{44},\qquad
c_{1063}=c_{1064}=0.
\end{gathered}
\]
Put
\[
\mu=A_c-2+\sum_{j=1}^{44}(A_c-3)A_c^j
       +2^{450}(2^{78}-1).
\]
Then the new directional coefficients are
\[
\begin{split}
D_C={}&\rho^{40}+\rho^{41}+\rho^{1100}+\rho^{2200}\\
 &+\sum_{e\in\mathcal E}c_e
       \bigl(\rho^{7631-e}+\rho^{8733-e}\bigr),\\
D_R={}&\rho.
\end{split}
\tag{7.4}
\]
The native masks are
\[
\begin{split}
M_F={}&4+\mu(\rho^{7631}+\rho^{8733})
 +\sum_{j=0}^9\rho^{1000+2j}
 +\sum_{j=0}^9\rho^{1021+2j}
 +\rho^{2195}+\rho^{3296},\\
M_C={}&B-1-\sum_{\substack{e\in\mathcal E\\e\ne1}}\rho^e
          -2\rho^{1063}.
\end{split}
\tag{7.5}
\]
Finally take the four ports exactly as in (2.1).
The code calls them \texttt{beta}, \texttt{K0}, \texttt{MC},
and \texttt{MF}; the stored payload for the last is
\texttt{MF\_source}.
The archive contains all four integers as deterministic compressed
binary payloads, their hashes, the transition table, the tile list,
and independent reconstruction code.
No variable exponent occurs in this definition: every sum is finite
over a displayed fixed range or the supplied finite tile list.

\subsection{Parity separates temporal and spatial tests}

All native copy positions are even. The temporal copy tests are
at the even positions \(1000+2j\), while the spatial copy tests
are at the odd positions \(1021+2j\).
This makes the native transport displacement \(1\) sufficient.
The factor \(BD_R\) in \(\kappa\) still moves to a neighboring
whole cell; only its displacement \emph{within} the cell is reduced.

At temporal target \(1000+2j\), the bare \(D_C\)-shift \(40\)
sees exactly the south copy \(960+2j\). The shift \(41\)
and the \(D_R\)-shift \(1\) would see unsupported odd positions.
The rotated native word contributes the north copy \(1000+2j\).
At spatial target \(1021+2j\), the \(D_C\)-shift \(41\)
sees the east copy \(980+2j\), and \(D_R\)'s shift \(1\)
sees the west copy \(1020+2j\). The shift \(40\) and the
rotated native word would require unsupported odd positions.
Selectors end at \(942\), and the odd dummy and anchors are
outside these tested ranges.

At anchor targets \(2195,3296\), the only bare-shift contributions
are \(1095+1100\) and \(1096+2200\), respectively.
The other bare shifts and \(D_R\) miss the native support.
Every reflected clause exponent is at least
\(7631-3296=4335\), above all these targets.
The twenty edge tests and two synchronization tests therefore
have exactly their inherited meanings.

\subsection{Clause isolation and alignment for arbitrary rotations}

Write \(\mathcal E=\mathcal S\cup\{2195,3296\}\), with
\(\max\mathcal S=1096\). The two clause bands differ by \(1102\).
No pair in \(\mathcal S\) has that difference, and
\[
2195-1102=1093\notin\mathcal E,\qquad
3296-1102=2194\notin\mathcal E.
\]
Since \(3296-2195=1101\), we have
\(1102\notin\mathcal E-\mathcal E\).
At either clause target the bare shifts and \(D_R\) also miss
\(\mathcal E\). Thus its coefficient is exactly
\(\sum_e c_e d_e\), with the same selector and copy equations
as the original language.

The upper-dummy spill position is \(1063+1=1064\), already in
\(\mathcal E\). Moreover \(L_c>2\max\mathcal E\).
Hence the two circular sets
\[
\{7631-e\bmod L_c:e\in\mathcal E\},
\qquad
\{8733-e\bmod L_c:e\in\mathcal E\}
\]
are disjoint. For every rotation, at least one clause band is clean.

For binary alignment, write a rotation as \(625t+\varepsilon\),
where \(0\le\varepsilon<625\).
The inherited carry-free bound implies that for
\(1\le\varepsilon\le623\) every rotated radix digit is even.
For \(\varepsilon=624\), only the single upper-dummy spill can
give an odd digit.
At an occupied marked cell, both tested anchors require odd
digits, excluding both kinds of misalignment.
For \(\varepsilon=0\), the sole ordered support pair at
difference \(1101\) is \((2195,3296)\):
\(2195-1101=1094\) is absent, whereas
\(3296-1101=2195\) is present.
There are no additional circular cases since
\(L_c>2\max\mathcal E\).
The anchor tests therefore force whole-cell alignment.
This checks arbitrary purported rotations; it does not assume
a desired stride divisibility in advance.

\subsection{Carry, population, origin, and canonical padding}

The values \(c_e\) are only relocated. Their mass and the
base-\(A_c\) clause inequalities remain unchanged.
For completeness the finite bounds checked in the archive are
\[
\begin{gathered}
\mathrm{mass}=(989+2)\left(2\sum_e c_e+6\right),\\
\mathrm{threshold}=\max(4\mathrm{mass}+8,\,2\mu+4,\,16)<\rho,\\
\frac{3\rho}{4}+\mathrm{mass}\le\rho-2,\qquad
\frac{\rho}{2}+1+\mathrm{mass}\le\rho-2.
\end{gathered}
\tag{7.6}
\]
The mass has \(461\) bits and the threshold \(529\) bits, versus
\(625\) bits of radix capacity.
The largest possible product exponents are bounded by
\[
8733+3296=12029<L_c,\qquad
1+3296=3297<L_c.
\]
There is no new intercell spill.
The masks satisfy
\[
\pc(M_F)=989,\quad \pc(M_C)=d_c-989,\quad
M_C\equiv2\pmod4,\quad M_F\equiv4\pmod8,
\]
and both native masks lie strictly between \(0\) and \(B-1\).
The low mask has its highest in-cell bit set, but neither mask
has a bit at or above the cell cutoff.

The minimum exponents of \(D_C,D_R\) are \(40,1\), both positive.
The low-origin unit therefore still comes only from the temporal
neighbor's selector \(0\). The canonical unique Start is not its
own temporal neighbor, preserving the special origin test.

Finally \(\rho\equiv2\pmod5\).
The paired clause contributions cancel modulo five, because their
exponent gap \(1102\) is \(2\pmod4\). The bare terms give
\[
D_C\equiv1+2+1+1=0\pmod5,\qquad
D_R\equiv2\pmod5,\qquad
2D_C-D_R\equiv3\pmod5.
\tag{7.7}
\]
Under the canonical power-of-five padding, let \(w_{\rm tile}\)
denote the tableau width (distinct from the core coordinate \(w\)).
The dummy coefficient controlling the actual packed rank is
\[
\Gamma=\rho^{1063}(q^2-1)
\bigl[1+q(D_C+BD_R+B^{w_{\rm tile}})\bigr]\equiv2\pmod5.
\tag{7.8}
\]
It is a unit modulo the required power of five.
Thus the inherited Boolean subset adjustment still realizes the
actual-rank stride congruence.
The unchanged tile language supplies independently large
power-of-five width and height padding. Together with the
preserved origin and strict slack bounds, this establishes
canonical completeness. The mask calculation (6.12) is unchanged,
so the simultaneous odd-\(i\), square-\(w\), odd-\(s\) refinements
also survive.

\begin{proposition}
The complete tuple (7.1)--(7.5) meets the native compiler's
soundness and canonical completeness interface. Consequently
it can be used in both Theorems~\ref{thm:positive} and
\ref{thm:integer}.
\end{proposition}
\begin{proof}
The edge, clause, occupancy, alignment, and carry arguments above
decode exactly the same Wang tilings as the inherited \(943\)-tile
language. Its Start/Forbidden and helical computation theorem
\cite{wang} decodes precisely blank halting computations.
Conversely the same language pads an actual finite computation;
(7.6)--(7.8) give the required packing, rank adjustment, and slack.
The native arithmetic compiler proof \cite{blank,ordinary}
therefore applies to the complete new tuple.
The supplied proof notes spell out the same adapter independently
of this compressed presentation.
\end{proof}

\subsection{Exact size reduction}

Since \(0<D_C<B\) and \(D_R=2^{625}\),
\[
\operatorname{bitlength}(\kappa)=d_c+625+1=9{,}766{,}251.
\]
This is the largest of the four actual port bit lengths.

\begin{table}[ht]
\centering\small
\begin{tabular}{lrrr}
\toprule
Measure & Previous Wang943 & New parity tuple & Reduction\\
\midrule
Largest port, bits & 13,201,251 & 9,766,251 & 3,435,000\\
Four raw payloads, bytes & 5,312,269 & 4,882,894 & 429,375\\
Four deterministic gzip files, bytes & 20,260 & 19,820 & 440\\
\bottomrule
\end{tabular}
\caption{The coefficient-height improvement is \(26.02\%\).
Compression lengths are separate storage measurements.}
\end{table}

This adapter alone does not reduce the tile count, variable count,
degree, or arithmetic operation count. It supplies a smaller exact
numeric backend for both new arithmetic constructions.

\section{Consistency equivalence, reproducibility, and limits}
\label{sec:verification}

\begin{theorem}[The requested equivalence]\label{thm:consistency}
Let the fixed four integers be the tuple in
Section~\ref{sec:coefficients}. ZFC proves
\[
\begin{split}
\neg\exists\mathbf x\in\Z_{>0}^{12}\,
 P_{12}(\mathbf x)=0
&\ \longleftrightarrow\ \Con(\mathrm{ZFC}),\\
\neg\exists\mathbf z\in\Z^{24}\,
 P_{24}(\mathbf z)=0
&\ \longleftrightarrow\ \Con(\mathrm{ZFC}).
\end{split}
\tag{8.1}
\]
The same equivalences hold for the shifted three-equation
systems (5.14), (6.14), in their stated domains, and for the
integer sum of squares (6.15).
\end{theorem}
\begin{proof}
Theorems~\ref{thm:positive} and \ref{thm:integer}, with the
coefficient adapter, give finite polynomial solutions if and
only if there is a finite blank halting computation of the
literal machine. The exact machine calibration identifies those
computations with finite ZFC refutations.
All reductions use finite coding, integer arithmetic, Pell
recurrences, finite algebraic extensions, and the indicated
classical number-theoretic facts; their proofs can be carried
out in ZFC.
Thus ZFC proves each existential statement equivalent to its
syntactic refutability statement. Negating gives (8.1).
The system shifts preserve existence in both directions,
and over the integers a sum of squares vanishes exactly when
every summand vanishes.
\end{proof}

This theorem does not decide which side holds.
In particular, under the usual consistency assumption, proving
one of these polynomials has no solutions within ZFC would
contradict the second incompleteness theorem.
The construction identifies the arithmetic statement whose
unprovability is at issue.

\subsection{What is literal and what is inherited}

The archive is intended as a reviewable repository contribution.
It contains:
\begin{itemize}
\item this \LaTeX{} source and its PDF, with the new proofs and
  exact nonvanishing degree certificates;
\item seven literal JSON arithmetic circuits: three for the
  positive construction and four for the integer construction;
\item Python generators and finite arithmetic diagnostics,
  plus an independent circuit/liveness/degree-bound checker;
\item the complete coefficient recipe, transition table, ordered
  Wang tile list, four compressed binary integers, and independent
  byte-for-byte reconstruction checks;
\item independent written reviews and the source copies needed
  to audit the reused norm, signed-domain, compiler, and endpoint
  interfaces;
\item source provenance, licenses, an artifact checksum manifest,
  and a root reproduction command.
\end{itemize}

The circuit files enumerate the supplied variables and the four
fixed ports separately, define every intermediate using a binary
operation, and identify their output equation or comparisons.
They can be expanded mechanically to ordinary polynomials,
although doing so loses useful sharing and greatly enlarges the
presentation. The degrees in Table~\ref{tab:new} are the degrees
of those expanded polynomials, not circuit-depth estimates.

The newly supplied proofs establish the two root reductions and
the coefficient adapter. The inherited native compiler,
Pell-index lemmas, Wang language theorem, and machine/refutation
calibration remain dependencies. Their pinned sources are
identified and relevant copies are included.
We do not claim to have independently reconstructed every
earlier compiler proof or rerun the complete formal endpoint.

\subsection{Reproduction and the scope of the checks}

From the archive root, with a standard Python 3 installation,
run
\begin{quote}\ttfamily
python3 reproduce.py
\end{quote}
No network, computer algebra package, or proof assistant is
required for the finite checks.
To rebuild the manuscript, use a standard \TeX{} installation:
\begin{quote}\ttfamily
pdflatex -interaction=nonstopmode -halt-on-error report.tex\\
pdflatex -interaction=nonstopmode -halt-on-error report.tex
\end{quote}

The arithmetic checks compare the compact norm evaluator with
an independent eight-dimensional conjugate-product evaluation
on signed assignments; exercise the inverse reconstruction
including negative strong roots; verify transport shifts;
audit live gates and port use; and check nonvanishing leading
forms modulo two primes.
The integer review separately derives the signed bounds,
the two boundary exclusions, and the exact weighted leading form.
The coefficient checks cover \(69{,}230\) basis convolution
identities, all-rotation support facts, carry/slack inequalities,
population counts, the five-adic unit, and the four exact payloads.
The inherited language checker tests \(425\) explicitly halting
control helices comprising \(33{,}800\) local tiles.

These tests detect transcription errors and verify the named
finite properties. Finite fixtures do not establish a universal
Diophantine equivalence. That conclusion depends on the proofs,
including the inherited compiler, rather than on a search for
zeros or a successful sample of computations.
No new Lean/Rocq proof object for (8.1) is supplied.

\subsection{Which presentation to use}

For twelve positive unknowns, (5.14) is the recommended system:
three equations, \(110\) operations, maximum degree \(1072\).
For a single equation with that domain, use (5.2).
For ordinary unrestricted integer unknowns, use (6.14) if a
system is acceptable and (6.4) if one equation is required.
The sum of squares is useful when its especially direct
zero-set argument outweighs its higher degree.

The contribution improves the inspected moderate-degree
twelve-positive and twenty-six-integer families.
It does not dominate the retained \(67\)-operation positive
polynomial, the \(59\)-operation larger-variable system,
the degree-\(172\) twenty-seven-integer product, or the
nine-natural-variable Jones specialization.
There is no proved global optimum among all Diophantine
representations of this consistency statement.

A further variable-count comparison is timely.
Sun's September 2026 preprint \cite{sun10}, Theorem 1.1,
states that every recursively enumerable subset of \(\N\)
has a representation using ten integer existential unknowns.
Its theorem gives a generic route to a ten-integer consistency
instance. We have not instantiated or independently audited
that compiler here, so this report makes no claim to a global
integer-variable record.

\subsection{Suggested next work}

The best immediate verification target is to formalize the
new signed separation and saturation composition as a small
extension of the existing core interface. The positive proof's
new obligation is similarly localized to (5.5)--(5.8).
For further size reduction, the most promising open comparison
is whether a weaker integer-valued or nonmonic inverse
construction can remove another guarded coordinate without
destroying soundness at zero or negative scales.
Any such proposal must recheck coherent root recovery,
the \(C=0\) boundary, exact leading cancellation, and simultaneous
canonical completeness. The present work stops at the improvements
proved and reproduced in this archive.


\begin{thebibliography}{99}\small
\bibitem{catalogue}
ProveIt, audited wave-3 reconciliation, snapshot
\texttt{a477d0fc8b6feda53796835564a988015e744d4d}.
\href{\repo\coll README.md}{Pinned collection README}.
\bibitem{endpoint}
ProveIt, \texttt{KEndpoint215.lean}, exact blank-halting/refutation endpoint.
\href{\repo Computability/BusyBeaver/Lean/BusyBeaver/ZFC/KEndpoint215.lean}{Pinned source}.
\bibitem{blank}
ProveIt, fixed blank-input \(67\)-operation compiler and helical proofs,
retained in package 21.
\href{\repo\coll 21-norm12/upstream/FIXED_BLANK_67_PROOF.md}{Pinned proof}.
\bibitem{ordinary}
ProveIt, ordinary strong norms and signed blank witnesses, retained in 21.
\href{\repo\coll 21-norm12/upstream/ORDINARY_BLANK_PROOF.md}{Pinned proof}.
\bibitem{signed}
ProveIt, normalized signed domain and canonical parity, retained in 21.
\href{\repo\coll 21-norm12/upstream/SIGNED_DOMAIN_29.md}{Pinned proof}.
\bibitem{root20}
ProveIt package 20, direct twelve-variable root-merger proof.
\href{\repo\coll 20-root12/amalgamation/ROOT12_PROOF.md}{Pinned proof}.
\bibitem{norm21}
ProveIt package 21, norm proof and economical evaluator.
\href{\repo\coll 21-norm12/proofs/NORM12_PROOF.md}{Pinned proof};
\href{\repo\coll 21-norm12/arithmetic/build_norm12.py}{pinned generator}.
\bibitem{monic22}
ProveIt package 22, monic merger and transport shift.
\href{\repo\coll 22-monic12/proofs/merger_proof.md}{Pinned merger};
\href{\repo\coll 22-monic12/proofs/RETAINED59_PROOF.md}{pinned shift}.
\bibitem{oddscale}
ProveIt package 23, odd signed half-binomial scale and exact population.
\href{\repo\coll 23-separating12/proofs/ODD_SCALE_27_PROOF.md}{Pinned proof}.
\bibitem{wang}
ProveIt package 19, Wang943 language and compact coefficient layout.
\href{\repo\coll 19-wang943/compiler/WANG943_COEFFICIENTS.md}{Pinned language proof};
\href{\repo\coll 19-wang943/compiler/WANG943_COMPACT_LAYOUT.md}{pinned layout}.
\bibitem{jones}
J.~P. Jones, \emph{Universal Diophantine equation},
Journal of Symbolic Logic \textbf{47} (1982), 549--571.
\href{https://doi.org/10.2307/2273588}{doi:10.2307/2273588}.
\bibitem{sun10}
Z.-W. Sun, \emph{Ten unknowns for Hilbert's tenth problem over the integers},
preprint, 1 September 2026, Theorem 1.1.
\href{https://arxiv.org/abs/2609.01594}{arXiv:2609.01594v1}.
\end{thebibliography}
\end{document}
